%% kapa: exam=1 round=36 year=2025 period=2 form=A subject=appraisal_law subject_ko=감정평가관계법규 seq=6 number=6
%% source: Q-Net 감정평가사 기출문제 (cst00309) artlSeq=5231525 file=raw/1008/5231525/2025년 제36회 감정평가사 1차시험 2교시 문제지 원본.pdf
%% conversion: kordoc 4.24.51 -> tools/gen1.py (auto); review=proofread-sonnet; flags=none
\begin{question}{6}
국토의 계획 및 이용에 관한 법령상 도시ㆍ군기본계획에 관한 설명으로 옳지 않은 것은?
\begin{choices}
  \item 광역도시계획이 수립되어 있는 지역에 대하여 수립하는 도시ㆍ군기본계획의 내용이 광역도시계획의 내용과 다를 때에는 도시ㆍ군기본계획의 내용이 우선한다.
  \item 도시ㆍ군기본계획에는 기후변화 대응 및 에너지절약에 관한 사항에 대한 정책 방향이 포함되어야 한다.
  \item 도시ㆍ군기본계획에서 정하는 생활권 간의 경계를 변경하는 생활권계획이 수립된 때에는 해당 계획이 수립된 생활권에 대해서는 도시ㆍ군기본계획이 수립된 것으로 본다.
  \item 도시ㆍ군기본계획 입안일부터 5년 이내에 토지적성평가를 실시한 경우에는 토지적성평가를 하지 아니할 수 있다.
  \item 수도권에 속하지 아니하고 광역시와 경계를 같이하지 아니한 시 또는 군으로서 인구 10만명 이하인 시 또는 군은 도시ㆍ군기본계획을 수립하지 아니할 수 있다.
\end{choices}
\end{question}
\answer{1}
